\section{Problems on the Binomial Theorem}
\subsection*{Question 1 \\ \small (JEE Mains 2024)}

Let $m$ and $n$ be the coefficients of the $7^{\text{th}}$ and 
$13^{\text{th}}$ terms respectively in the expansion of 
\[
\left(\frac{1}{3}x^{\frac{1}{3}} + \frac{1}{2x^{\frac{2}{3}}}\right)^{18}.
\]
Find $\left(\frac{n}{m}\right)^{\frac{1}{3}}$.

\subsection*{Question 2 \\ \small(JEE Mains 2022)}

If the sum of the coefficients of all positive even power of $x$ in the binomial expansion of $(2x^3+ \frac{x}{3})^{10}$ is $5^{10} - \beta.3^9$, then $\beta$ is equal to

\subsection*{Question 3 \\ \small(JEE Adv 2013, JEE Mains 2020)}
Coefficient of three consecutive term in the expansion of $(1+x)^{n+5}$ are in ratio $5:10:14$ then n is equal to?

\subsection*{Question 4 \\ \small(JEE Mains 2024)}
Let the coefficient of $x^r$ in the expansion of \\
\[
(x+3)^{n-1}+(x+3)^{n-2}(x+2)+(x+3)^{n-3}(x+2)^2+(x+3)^{n-4}(x+2)^3+...+(x+2)^{n-1}
\]
be $a_r$. If $\sum_{r=0}^{n}a_r = \beta^n - \gamma^n, \beta,\gamma \epsilon N$, then the value of $\beta^n+\gamma^n = $? 
\newpage
\subsection*{Question 5 \\ \small(JEE Mains 2022)}
If the coefficient of $x^{10}$ in the binomial expansion of \[
({\frac{\sqrt{x}}{5^{\frac{1}{4}}}+{\frac{\sqrt{5}}{x^{\frac{1}{3}}}}})
\] is $5^kl$, where $l,k \epsilon N$ and $l$ is co-prime to 5, then k equal to $=$?

\subsection*{Question 6}
The coefficient of $x^{301}$ in \[
(1+x)^{500}+x(1+x){499}+x^{2}(1+x)^{498}+...+x^{500}
\] is?

\subsection*{Question 7}
The number of integral terms in the expansion of \[
(7^{(\frac{1}{2})}+11^{(\frac{1}{6})})^{824}
\]

\subsection*{Question 8 \\ \small(JEE Mains 2019)}
The total number of irrational terms in the binomial expansion of \[
(7^{\frac{1}{5}}-3^{\frac{1}{10}})^{60}
\]

\subsection*{Question 9 \\ \small(JEE Mains 2022)}
The term independent of x in the expression of \[
(1-x^2-3x^3)({\frac{5}{2}x^3-\frac{1}{5x^2}})^{11}
\], $x \ne 0$

\subsection*{Question 10 \\ \small(JEE Mains 2023)}
If the coefficient of $x$ and $x^2$ in $(1+x)^p(1-x)^q$ are 4 and -5 respectively, then $2p+3q$ equal to?

\subsection*{Question 11 \\ \small(JEE Mains 2023)}
If the constant term in the expansion of given expression is $p$, then $108p$ equals to? \[
(1+2x-3x^2)(\frac{3}{2}x^2-\frac{1}{3x})^9
\]


\subsection*{Question 12 \\ \small(JEE Mains 2024)}
In the expansion of \[
(1+x)(1-x^2)(1+\frac{3}{x}+\frac{3}{x^2}+\frac{1}{x^3})^5, x\ne 0,
\] then sum of the coefficient of $x^3$ and $x(-13)$ is equal to?

\subsection*{Question 13 \\ \small(JEE Mains 2024)}
If $A$ denotes the sum of all the coefficient in the expansion of $(1=3x+10x^{2})^n$ and $B$ denotes the sum of all the coefficient in the expansion of $(1+x^2)$, then:

\begin{enumerate}
    \item $B=A^3$
    \item $3A=B$
    \item $A=3B$
    \item $A=B^3$
\end{enumerate}

\subsection*{Question 14 \\ \small(JEE Mains 2019)}
The coefficient of $x^{18}$ in the product of \[
(1+x)(1-x)^{10}(1+x+x^2)^9
\]

\subsection*{Question 15 \\ \small(JEE Mains 2021)}
Let $n$ be the positive integer. Let
\[
A = \sum_{k=0}^{n}(-1)^k[(\frac{1}{2})^k+(\frac{3}{4})^k+(\frac{7}{8})^k+(\frac{15}{16})^k+(\frac{31}{32})^k]
\] If $63A = 1-\frac{1}{2^{30}}$ then $n$ is equal to?

\subsection*{Question 16 \\ \small(JEE Mains 2022)}
Let for the $9^{th}$ term in the binomial expansion of $(3+6x)^n$, in the increasing power of $6x$, to be the greatest for $x=\frac{3}{2}$, the least value of $n$ is $n_o$. If $k$ is the ration of the coefficient of $x^6$ to the coefficient of $x^3$, then $k+n_o$ is equal to.

\subsection*{Question 17 \\ \small(JEE Mains 2024)}
\[
{}^{n-1}C_{r} = (k^2-8){}^{n}C_{r+1}
\]
If and only if: 
\begin{enumerate}
    \item $2\sqrt{2} < k \leq 3$
    \item $2\sqrt{3}< k \leq 3\sqrt{2}$
    \item $2\sqrt{3}< k< 3\sqrt{3}$
    \item $2\sqrt{2}k<k<2\sqrt{3}$
\end{enumerate}

\subsection*{Question 18 \\ \small(JEE Mains 2020)}
The value of $\sum_{r=0}^{20} {}^{50-r}C_{r+6}$ is equal to: 
\begin{enumerate}
    \item ${}^{50}C_{6} - {}^{30}C_{6}$
    \item ${}^{51}C_{7}+{}^{30}C_{7}$
    \item ${}^{51}C_{7} - {}^{30}C_{7}$
    \item ${}^{50}C_{7} - {}^{30}C_{7}$
\end{enumerate}

\subsection*{Question 19 \\ \small(JEE Mains 2021)}
Let ${}^{n}C_{r}$ denote the binomial coefficient of $x^r$ in the expansion of $(1+x)^n$. If \[
\sum_{k=0}^{10} (2^k+3k){}^{10}C_{k} = \alpha .3^{10}+\beta.2^{10}, \alpha, \beta \in R
\]
Then $\alpha+\beta=$?
\newpage

\subsection*{Question 20 \\ \small(JEE Mains 2023)}
Suppose \[
\sum_{r=0}^{2023}r^2.{}^{2023}C_{r} = 2023 \times\alpha\times2^{2022}.
\]
Then the value of $\alpha$ is.

\subsection*{Question 21 \\ \small(JEE Mains 2023)}
If, \[
\frac{1}{n+1}{}^{n}C_{n}+\frac{1}{n}{}^{n}C_{n-1}+\cdots+{}^{n}C_{0} = \frac{1023}{10}
\]. Then n equal to.

\subsection*{Question 22 \\ \small(JEE Mains 2024)}
Let \[
\alpha = \sum_{k=0}^{n} (\frac{({}^{n}C_{k})^2}{k+1})
\]
\[
\beta = \sum_{k=0}^{n-1}(\frac{({}^{n}C_{k})({}^{n}C_{k+1})}{k+2}) 
\]
If $5\alpha = 6\beta$, then n equals.

\subsection*{Question 23 \\ \small(JEE Mains 2019)}
The value of $r$ for which \[
{}^{20}C_{r}\cdot{}^{20}C_{0}+{}^{20}C_{r-1}\cdot{}^{20}C_{1}+{}^{20}C_{r-2}\cdot{}^{20}C_{2}+\cdots+{}^{20}C_{0}\cdot{}^{20}C_{r}
\] is maximum

\subsection*{Question 24 \\ \small(JEE Adv. 2018)}
Let , \[
X = ({}^{10}C_{1})^2+2({}^{10}C_{2})^2+\cdots+10({}^{10}C_{10})^2
\]
Where $r \in \{1,2,3.\cdots,10\} $ and ${}^{10}C_{r}$ denote binomial coefficient. Then value of $\frac{1}{1430}X$ is?

\subsection*{Question 25 \\ \small(JEE Mains. 2020)}
Let \[
(2X^2+3x+4)^{10} = \sum^{20}_{r=0}a_rx^r.
\]
Then $\frac{a_7}{a_{13}}$ is equal to?
\subsection*{Question 17}
\[ 
P = \sum_{j=0}^{n}\sum^{n}_{i=0}({}^{n}C_{i}+{}^{n}C_{j})
\]
\[
Q = \sum_{j=0}^{n}\sum^{n}_{i=0}({}^{n}C_{i}+{}^{n}C_{j}) 
\]
Where $i=j$
\[
R = {\sum^{n}_{0\leq i<j \leq n}\sum^{n}}_{}({}^{n}C_{i}+{}^{n}C_{j})
\]
\[
S = {\sum^{n}_{0\leq i \leq j \leq n}\sum^{n}}_{}({}^{n}C_{i}+{}^{n}C_{j})
\]

\subsection*{Question 26 \\ \small(JEE Mains. 2022)}
\[
\sum_{i,j=0}^{n} {}^{n}C_{i} \cdot {}^{n}C_{j} 
\] is equal to
\begin{enumerate}
    \item $2^{2n}-{}^{2n}C_{n}$
    \item $2^{2n-1}-{}^{2n-1}C_{n-1}$
    \item $2^{2n} - \frac{1}{2} \cdot{}^{2n-1}C_{n-1}$
    \item $2^{2n} + {}^{2n-1}C_{n-1}$
\end{enumerate}

\subsection*{Question 27}
In the expansion if $(1+x+x^2)^5$ find:
\begin{enumerate}
    \item Coefficient of $x^4$
    \item Number of terms
\end{enumerate}

\subsection*{Question 28 \\ \small(JEE Mains. 2021)}
If the coefficient of $a^7b^8$ in the expansion of $(a+2b+4ab)^{10}$ is $K.2^{16}$, then $K$ is equal to?

\subsection*{Question 29 \\ \small(JEE Mains. 2023)}
Let,
\[
(a+b+cx^2) = \sum_{i=0}^{20}p_ix^i
\]
\[
a,b,c \in N
\].
If $p_1 = 20$ and $p_2 = 210$ then $2(a+b+c)$ is equal to

\subsection*{Question 30 \\ \small(JEE Mains. 2023)}
The coefficient of $x^7$ in $(1-x+2x^3)^{10}$ is


\subsection*{Question 31}
Find the remainder for the following.
\begin{enumerate}
    \item $7^{2022}+3^{2022}$ divided by 5.
    \item $\{\frac{2^{403}}{15}\}$, Where \{.\} $\equiv$
    \item $\{\frac{2^{200}}{8}\}$, Where \{.\} $\equiv$
    \item $(11)^{1011} + (1011)^{11}$ is divided by 9
\end{enumerate}

\subsection*{Question 32 \\ \small(JEE Mains. 2024)}
Remainder when $64^{32^{32}}$ is divided by 9 is equal to.

\subsection*{Question 33 \\ \small(JEE Mains. 2024)}
Remainder when $428^{2024}$ is divided by 21 is equal to.

\subsection*{Question 34 \\ \small(JEE Mains. 2022)}
Remainder when $2021^{2023}$ is divided by 7 is equal to.

\subsection*{Question 35}
Find the last digit of $3^{104}$

\subsection*{Question 36}
If $(5+2\sqrt{6})^n = I+f$, where $I \in N, n \in N$ and $0 \leq f< 1$, then prove that $I=\frac{1}{1-f}-f$.

\subsection*{Question 37 \\ \small(JEE Mains. 2023)}
Let $x = (8\sqrt{3}+13)^{13}$ and $y = (7\sqrt{2}+9)^3$. If [t] denotes the greatest integer $\leq t$, then: 
\begin{enumerate}
    \item $[X]+[y]$ is even
    \item $[x]$ is odd but $[y]$ is even
    \item $[x]$ is even but $[y]$ is odd
    \item $[x]$ and $[y]$ are both odd
\end{enumerate}
\subsection*{Question 38 \\ \small(JEE Mains. 2019)}
The coefficient of $t^4$ in the expansion of $\left( \frac{1-t^6}{1-t} \right)^3$ is?